\chapter{Monitoring and Retraining}\label{ch:ml:monitoring-and-retraining}

On a Tuesday the model stopped trading: every prediction was inside the threshold. Nothing had failed. An upstream vendor had
begun sending volumes in lots instead of shares, and the model was obediently scaling them. No test had broken, because every
number was still a number; the only sign was that the desk had nothing to do. This chapter builds the monitors that would have
paged someone before lunch, and asks of each the question that decides whether it is worth having: how long does it take to
see a real failure when it is only allowed to cry wolf once a month? Three failures are planted in a simulated production run: a
unit change upstream, a reversal of the strongest feature's effect, and a feature whose effect fades while its spread grows.
Monitors on the inputs see the first on the day it happens and are blind to the second; the outcome monitor sees the second in a
day and cannot see the first for months; nothing sees the third's lost accuracy, and the input monitors catch it only through its
side effect, six weeks in.

\section{What to monitor: inputs, outputs, outcomes}

\begin{definition}[Model monitoring, prediction drift]\label{def:ml:monitoring-and-retraining:monitoring}
\emph{Model monitoring}\index{model monitoring} is the continuous comparison of a production model's inputs, outputs and
realised outcomes with what was expected when it was validated, with alarms that page a person when they differ by more than
chance allows. \emph{Prediction drift}\index{prediction drift} is a change in the distribution of the model's outputs relative to
the validation period, whatever its cause.
\end{definition}

Chapter~12 separated \omterm{def:ml:online-learning-and-drift:drift}{covariate shift}, a change in the inputs, from \omterm{def:ml:online-learning-and-drift:drift}{concept drift}, a change in how the inputs relate to the
target, and built \omterm{def:ml:online-learning-and-drift:detector}{drift detectors} on a model's errors. A production model is watched at three layers, because each sees failures
the others miss. Inputs are available at once and in bulk: every feature of every name, every day, compared with the training
reference, after the schema and delivery checks of chapter~15 and the freshness checks of chapter~24. Outputs are available at
once too: the distribution of predictions, the share of names whose forecast clears the trading threshold, the positions and
turnover they imply. Outcomes (the realised information coefficient, the P\&L) arrive only after the forecast horizon, and they
are noisy: a model with a true daily information coefficient of 0.10 shows a standard deviation of about 0.07 from one day to the
next. Outcomes are what matter, and they are the slowest to speak.

The chapter's production model is deliberately plain, so that every failure is known exactly. Each day 200 names carry five
standardised features; the next day's returns load on the first three with coefficients 0.08, $-0.06$ and 0.05, plus noise of
unit variance. A linear model is fitted on 60 days and then runs for 400 more; it trades the names whose forecast is larger in
absolute value than the training median, half of them. (Chapter~29's order-book model is monitored with the same code at the end
of the book.) Three failures start on day 200, each in 20 independent runs (\cref{lst:ml:mon:day}):
\begin{itemize}
\item \emph{unit change}: the second feature arrives a hundred times smaller, as volumes in lots would;
\item \emph{reversal}: the first feature's effect changes sign, a regime change the inputs do not show;
\item \emph{slow fade}: over 150 days the third feature's effect falls to zero while its spread grows by half.
\end{itemize}

\section{Drift statistics and alarm design}

\begin{definition}[Population stability index]\label{def:ml:monitoring-and-retraining:psi}
Cut a reference sample of a variable into $B$ bins (here its deciles) with shares $r_1,\dots,r_B$, and let $q_1,\dots,q_B$ be the
shares of a new sample in the same bins. The \emph{population stability index}\index{population stability index} is
\[
\mathrm{PSI}=\sum_{b=1}^{B}(q_b-r_b)\ln\frac{q_b}{r_b},
\]
the symmetrised Kullback--Leibler divergence between the two binned distributions (empty bins are floored at a small share).
\end{definition}

The index comes from credit scoring, where the rules of thumb are that below 0.10 a population has shifted little and above 0.25
it has shifted significantly; Yurdakul and Naranjo point out that these benchmarks are used without reference to their error
rates. The Kolmogorov--Smirnov statistic of Book~4 (chapter~12), the largest gap between the two empirical distribution
functions, is the other workhorse (\cref{lst:ml:mon:psi}). Both are computed each day on each feature against the training
reference, and the monitor takes the largest over the five features. With many features, one test per feature multiplies the
false alarms; Rabanser and co-authors, comparing ways of detecting dataset shift, found two-sample tests on a reduced
representation of the inputs the most effective: a single multivariate test, like chapter~16's \omterm{def:ml:generative-models-and-synthetic-data:tests}{classifier two-sample test}. The output monitors compute the index on the predictions
and track the share of names inside the trading threshold, which is where the Tuesday's failure would show first.

An alarm is a threshold, and a threshold is a false-alarm rate. The chapter fixes one false page a month, 1 in 21 trading days,
and sets each threshold on 20 runs of 400 days without failures (\cref{lst:ml:mon:calib}). A page is counted when an alarm
starts, not for every day it lasts, because a statistic that stays above its threshold for a week pages the on-call person once.
The calibrated thresholds are 0.115 for the largest feature index, 0.115 for the largest Kolmogorov--Smirnov statistic, 0.089
for the index of the predictions and 0.07 for the change in the share inside the threshold. The first is above the credit-scoring
rule of thumb: with 200 names a day, deciles hold 20 names each, and sampling noise alone pushes the largest of five indices past
0.10 on more than one day in twenty. A threshold copied from a rule of thumb would page more often than the team agreed to, and
a team paged too often stops reading the pages.

Two numbers describe each monitor's response to a failure. The \emph{first alarm} is the number of days from the failure to the
monitor's first page; since every monitor pages falsely once a month, a monitor that sees nothing still pages after about two
weeks, and a first alarm is only a detection when it comes well before that. The \emph{sustained alarm} is the number of days to
the first 21-day window with at least five alarm days, a page that keeps coming back; without a failure, no monitor's median
reaches one within 200 days.

\section{Performance alarms on noisy outcomes}

\begin{definition}[Performance alarm]\label{def:ml:monitoring-and-retraining:performance-alarm}
A \emph{performance alarm}\index{performance alarm} is a monitor on realised outcomes (the information coefficient, hit rate,
P\&L or mark-outs of the model's trades) that pages when their level falls below what validation led the firm to expect, by more
than the outcomes' own noise allows at the chosen false-alarm rate.
\end{definition}

The chapter's \omterm{def:ml:monitoring-and-retraining:performance-alarm}{performance alarm} is the one-sided CUSUM test of Book~7 (chapter~13) on the daily information coefficient
(\cref{lst:ml:mon:cusum}): it accumulates the shortfall of each day's coefficient below its expected level of 0.103, less an
allowance $k$ of half its daily standard deviation (0.035), and pages when the sum exceeds $h$, then starts again. At one false
page a month, $h$ is 0.11, about one and a half days' standard deviation: the monitor can afford to wait very little.

That tells how much it can see. A fall of the mean coefficient by $\Delta$, measured with daily noise $\sigma$, needs about
\[
n\approx\Bigl(\frac{(z_{0.95}+z_{0.90})\,\sigma}{\Delta}\Bigr)^2
\]
days to be told from noise with a one-sided test at 5\% and 90\% power. After the reversal the coefficient falls from 0.103 to
0.003, and $n$ is about 4 days; after the unit change it falls to 0.087, and $n$ is about 150 days; at the end of the slow fade it
has fallen to 0.081, and $n$ is about 80 days, but the fade takes 150 days to get there. Outcomes cannot see small losses of
accuracy on any horizon a desk would accept; they see large ones fast.

\begin{table}[htbp]
\centering\small
\begin{tabular}{@{}lccccccc@{}}
\toprule
 & none & \multicolumn{2}{c}{unit change} & \multicolumn{2}{c}{reversal} & \multicolumn{2}{c}{slow fade}\\
\cmidrule(lr){2-2}\cmidrule(lr){3-4}\cmidrule(lr){5-6}\cmidrule(lr){7-8}
monitor & first & first & sustained & first & sustained & first & sustained\\
\midrule
largest feature PSI & 12 & 0 & 0 & 12 & never & 12 & 42\\
largest feature KS & 10 & 0 & 0 & 10 & never & 8 & 60\\
prediction PSI & 17 & 0 & 0 & 17 & never & 19.5 & never\\
share inside threshold & 11 & 0 & 0 & 11 & never & 9.5 & never\\
IC CUSUM & 13 & 9.5 & never & 1 & 0 & 13 & never\\
\bottomrule
\end{tabular}
\caption{Days from the failure (day 200) to each monitor's first alarm and to a sustained alarm (five alarm days in 21), medians
over 20 runs, every monitor calibrated to one false page a month; ``none'' is the same count on runs without a failure, what a
blind monitor scores; ``never'' means no sustained alarm in the 200 days that follow. Data: \texttt{ml\_monitor.delays}.}
\label{tab:ml:mon:delays}
\end{table}

\Cref{tab:ml:mon:delays} is the chapter's result. The unit change is seen by every input and output monitor on its first day:
the second feature's index is about 7, sixty times its threshold, and the share of names inside the trading threshold rises from
50\% to 63\%, a quarter fewer names traded. The outcome monitor pages once after a median of 9.5 days and never insists. The
reversal is invisible to every input and output monitor (their columns equal the ``none'' column: the inputs and the
distribution of predictions have not changed at all), and the CUSUM pages within two days in 13 of the 20 runs and keeps paging. The
slow fade's lost accuracy is seen by nothing; its growing spread is seen by the input monitors, sustained after a median of 42
days for the index and 60 for the Kolmogorov--Smirnov statistic (\cref{fig:ml:mon:fade}). Had the fade come without the change
in spread, nothing in the table would have caught it.

\begin{omfigure}
\begin{tikzpicture}
\begin{axis}[width=11.5cm,height=6cm,xlabel={production day},ylabel={alarm days per month},xmin=0,xmax=400,ymin=0,ymax=22,
  tick label style={font=\small},label style={font=\small},grid=major,grid style={gray!25},legend pos=north west,
  legend cell align=left,legend style={font=\small},table/col sep=comma]
\addplot[omProp,very thick,mark=*,mark size=1.3pt] table[x=day,y=psi]{figdata/ml/27-monitoring-and-retraining/fade.csv};
\addplot[omDef,very thick,dashed,mark=square*,mark size=1.3pt,mark options={solid}] table[x=day,y=ks]{figdata/ml/27-monitoring-and-retraining/fade.csv};
\addplot[gray,thick,mark=triangle*,mark size=1.5pt] table[x=day,y=predpsi]{figdata/ml/27-monitoring-and-retraining/fade.csv};
\addplot[omThm,very thick,dotted,mark=diamond*,mark size=1.6pt,mark options={solid}] table[x=day,y=cusum]{figdata/ml/27-monitoring-and-retraining/fade.csv};
\addplot[black,thin,dashed,forget plot] coordinates {(200,0) (200,22)};
\legend{largest feature PSI,largest feature KS,prediction PSI,IC CUSUM}
\end{axis}
\end{tikzpicture}

\omcaption{Alarm days per month (21 days) under the slow fade, in 25-day windows, averaged over 20 runs; the fade starts on day
200 (dashed) and ends on day 350. A blind monitor alarms on about one day a month. Data:
\texttt{ml\_monitor.alarm\_rates}.}\label{fig:ml:mon:fade}
\end{omfigure}

\begin{method}[Designing a model's monitors]\label{met:ml:mon:design}
\begin{enumerate}
\item Monitor inputs, outputs and outcomes: each layer sees failures the others cannot.
\item Choose the false-page rate the team will actually answer, and calibrate every threshold to it on history or simulation
without failures; do not copy thresholds from rules of thumb.
\item Count pages at the onset of an alarm, and distinguish a first page from a sustained alarm.
\item Measure each monitor's delay on planted failures, and compare it with a blind monitor's; a monitor that never beats that
column is noise.
\item Accept that outcome monitors see only large losses quickly; cover small ones with scheduled revalidation.
\end{enumerate}
\end{method}

\section{Shadow, canary and champion--challenger rollouts}

\begin{definition}[Shadow deployment, champion--challenger]\label{def:ml:monitoring-and-retraining:shadow}
In a \emph{shadow deployment}\index{shadow deployment} a new model receives the production inputs and records its predictions,
but its decisions are not executed; its outcomes are computed as if they had been. \emph{Champion--challenger}\index{champion--challenger}
is the discipline of keeping the production model (the champion) in place until a challenger, run beside it on the same data,
has shown with a test chosen in advance that it is better.
\end{definition}

A shadow run costs nothing but computation and gives a clean comparison, because champion and challenger see the same names on
the same days. After the reversal, the CUSUM pages within days. A challenger is refitted on the 40 days that
follow (its coefficient on the first feature is $-0.081$ on average over the runs, against the true $-0.08$) and runs in shadow
from day 240. Its mean information coefficient is 0.105, the champion's 0.003. The comparison uses Book~7's always-valid
sequential test (chapter~21) on the daily differences (\cref{lst:ml:mon:shadow}), which may be looked at every day without
inflating its error rate; the standard error is estimated from the days seen, so the test gives no verdict for ten days, since on
three days a small sample standard deviation makes any difference look certain. The challenger is promoted in all 20 runs, after
a median of 10 shadow days and at most 19 (\cref{fig:ml:mon:shadow}). Without a failure, a challenger refitted on 40 ordinary
days is promoted in 2 of 100 runs, within the test's 5\%.

\begin{omfigure}
\begin{tikzpicture}
\begin{axis}[width=11cm,height=5.8cm,xlabel={shadow day},ylabel={always-valid p-value},ymode=log,xmin=1,xmax=20,
  ymin=1e-7,ymax=2,tick label style={font=\small},label style={font=\small},grid=major,grid style={gray!25},
  table/col sep=comma]
\addplot[omProp,thick,mark=*,mark size=1.2pt] table[x=day,y=run0]{figdata/ml/27-monitoring-and-retraining/shadow.csv};
\addplot[omDef,thick,mark=square*,mark size=1.2pt] table[x=day,y=run1]{figdata/ml/27-monitoring-and-retraining/shadow.csv};
\addplot[omThm,thick,mark=triangle*,mark size=1.4pt] table[x=day,y=run2]{figdata/ml/27-monitoring-and-retraining/shadow.csv};
\addplot[gray,thick,mark=diamond*,mark size=1.4pt] table[x=day,y=run3]{figdata/ml/27-monitoring-and-retraining/shadow.csv};
\addplot[black,thick,mark=o,mark size=1.2pt] table[x=day,y=run4]{figdata/ml/27-monitoring-and-retraining/shadow.csv};
\addplot[black,thin,dashed] coordinates {(1,0.05) (20,0.05)};
\end{axis}
\end{tikzpicture}

\omcaption{Always-valid p-values of the challenger's advantage over the champion, day by day in shadow after the reversal, five
of the 20 runs; the test gives no verdict for ten days; dashed: 0.05, where the challenger is promoted. Data:
\texttt{ml\_monitor.shadow\_run}.}\label{fig:ml:mon:shadow}
\end{omfigure}

Shadow comparisons measure forecasts, not the market's reaction to trading on them; a model whose orders move prices is compared
in a canary deployment (Book~7, chapter~21), trading a small share of capital. The canary's value is in what it limits. The
chapter promotes a challenger with a sign error on its second feature, a bug a shadow run would have caught but suppose it did
not: the CUSUM rolls it back after a median of 3 days (at most 18). The information coefficient given up, summed over those days
and weighted by capital, is 0.26 IC-days when the challenger trades the whole book and 0.026 when it trades a tenth of it. The
rollback itself is a registry action (\cref{lst:ml:mon:rollback}): the version that was in production before is promoted back,
with the reason, into chapter~25's \omterm{def:ml:experiment-tracking-and-reproducibility:registry}{audit trail}.

\section{Retrain, roll back, or retire}

Each alarm calls for a different action, and the wrong one hides the failure. An input alarm with no fall in outcomes (the unit
change) calls for fixing the data: retraining on volumes in lots would teach the model a new scale for one feature and bury the
vendor's error in its coefficients, until the vendor corrected it. An outcome alarm with calm inputs (the reversal) says the
relation has changed: cut the model's size or stop it (the kill switch of Book~11, chapter~27, for a model trading at speed),
refit on data from after the change once there is enough, and promote through shadow and canary. A slow loss that no monitor
sees is the case for scheduled retraining (chapter~12's \omterm{def:ml:online-learning-and-drift:schedule}{retraining schedule}) and periodic revalidation of the model against a
simple benchmark. A model is retired when its refitted challengers no longer beat a null: when the effect is gone, retraining
only fits noise faster. Sculley and co-authors list ``changes in the external world'' among the debts of machine-learning
systems, and Breck and co-authors make monitoring half of their rubric for a model's production readiness. In the interviews
of machine-learning engineers by Shankar and co-authors, false-positive alerts were the most commonly discussed pain point, and
the alert fatigue that follows leads teams to silence alerts and miss real drops: the reason to count pages.

\begin{method}[After an alarm]\label{met:ml:mon:after}
\begin{enumerate}
\item Input or output alarm: check the data path first (units, schema, vendor changes); fix the data, not the model.
\item Outcome alarm with calm inputs: reduce or stop trading, refit on post-change data, and promote only through shadow and a
canary.
\item A challenger is promoted by a sequential test chosen in advance, and rolled back by the \omterm{def:ml:monitoring-and-retraining:performance-alarm}{performance alarm}, both recorded in
the registry.
\item Retrain on a schedule for the losses no monitor sees; retire the model when retraining no longer beats a null.
\end{enumerate}
\end{method}

\section{Tutorial: the Tuesday the model went quiet}\label{tut:ml:monitoring-and-retraining:tut}

\begin{tutorial}
\textbf{Goal.} Plant three failures in a simulated production run, calibrate five monitors to one false page a month, measure
their delays, promote a challenger from shadow and roll back a faulty one. \textbf{End state:} \cref{tab:ml:mon:delays},
\cref{fig:ml:mon:fade}, \cref{fig:ml:mon:shadow}.
\begin{tutsteps}
\item \textbf{The three failures.}
\omcode{code/ml/27-monitoring-and-retraining/python/ml_monitor.py}{44}{58}{One day of production data, with the planted failures.}\label{lst:ml:mon:day}
\item \textbf{Drift statistics.}
\omcode{code/firm/mlmonitor/firm_mlmonitor.py}{49}{65}{The population stability index and the Kolmogorov--Smirnov statistic against a prepared reference.}\label{lst:ml:mon:psi}
\item \textbf{Thresholds at a page rate.}
\omcode{code/firm/mlmonitor/firm_mlmonitor.py}{94}{110}{Alarm onsets, and the threshold that pages at a chosen rate.}\label{lst:ml:mon:calib}
\item \textbf{The \omterm{def:ml:monitoring-and-retraining:performance-alarm}{performance alarm}.}
\omcode{code/firm/mlmonitor/firm_mlmonitor.py}{75}{86}{A one-sided CUSUM on the daily information coefficient.}\label{lst:ml:mon:cusum}
\item \textbf{The shadow test.}
\omcode{code/firm/mlmonitor/firm_mlmonitor.py}{118}{127}{Always-valid p-values of a challenger's advantage, with no verdict before ten days.}\label{lst:ml:mon:shadow}
\item \textbf{Rollback.}
\omcode{code/firm/mlmonitor/firm_mlmonitor.py}{130}{143}{Putting the previous production version back through the registry.}\label{lst:ml:mon:rollback}
\item \textbf{Run} \texttt{ml\_monitor.thresholds()}, \texttt{delays()}, \texttt{shadow()}, \texttt{rollbacks()} and
\texttt{fig\_monitor.py} (about 15 seconds on one core).
\end{tutsteps}
\textbf{What to change next.} Pool the last 20 days of each feature into one index and calibrate it the same way: it sees the
slow fade's spread sooner, but its false alarms come in spells, and one page a month means something different for it
(\cref{exo:ml:monitoring-and-retraining:7}).
\end{tutorial}

\section{Build: model monitoring}\label{bld:ml:monitoring-and-retraining:mlmonitor}

\begin{build}
\bfield{Purpose} Monitors on a production model's inputs, outputs and outcomes, calibrated to a false-page rate, and the
promotion and rollback decisions they feed.
\bfield{Interface} \texttt{Reference(ref, bins)}; \texttt{psi(ref, x)}; \texttt{ks(ref, x)}; \texttt{daily\_ic(pred, y)};
\texttt{Cusum(target, k, h).update(v)}; \texttt{calibrate\_daily(null, rate)}; \texttt{onsets(alarms)};
\texttt{calibrate\_pages(null, rate)}; \texttt{first\_alarm(stats, threshold, start)}; \texttt{shadow\_verdict(ic\_challenger,
ic\_champion, tau, min\_days)} on \texttt{firm.abtest}; \texttt{rollback(registry, name, ts, reason)} on \texttt{firm.exptrack}.
\bfield{Rules} Every threshold is calibrated on data without failures at a stated rate; pages are counted at onsets; promotions
use a test fixed in advance; rollbacks go through the registry with a reason.
\bfield{Acceptance tests} \texttt{code/firm/mlmonitor/tests/}: the index is zero on its own reference and grows with a shift; the
Kolmogorov--Smirnov statistic and the information coefficient match SciPy, ties included; the CUSUM pages on a drop and resets;
calibrated thresholds page at their rate; the shadow test gives no verdict before ten days and promotes a null challenger in at
most 5\% of runs; a rollback restores the previous version and keeps the \omterm{def:ml:experiment-tracking-and-reproducibility:registry}{audit trail} intact.
\bfield{Stretch} Multivariate drift (a \omterm{def:ml:generative-models-and-synthetic-data:tests}{classifier two-sample test}, chapter~16) and monitors on P\&L and mark-outs, calibrated the
same way.
\end{build}

\begin{omsources}
\item S.~Rabanser, S.~Günnemann and Z.~C.~Lipton, ``Failing loudly: an empirical study of methods for detecting dataset
shift'', \emph{NeurIPS}, 2019.
\item E.~Breck, S.~Cai, E.~Nielsen, M.~Salib and D.~Sculley, ``The ML test score: a rubric for ML production readiness and
technical debt reduction'', \emph{IEEE International Conference on Big Data}, 2017.
\item B.~Yurdakul and J.~D.~Naranjo, ``Statistical properties of the population stability index'', \emph{Journal of Risk Model
Validation}, 2020.
\item D.~Sculley and co-authors, ``Hidden technical debt in machine learning systems'', \emph{NeurIPS}, 2015.
\item S.~Shankar, R.~Garcia, J.~M.~Hellerstein and A.~G.~Parameswaran, ``Operationalizing machine learning: an interview
study'', arXiv:2209.09125, 2022; published in \emph{Proceedings of the ACM on Human-Computer Interaction}, 2024.
\end{omsources}

\section{Exercises}

\begin{exercise}[$\star$]\label{exo:ml:monitoring-and-retraining:1}
For each of the three layers (inputs, outputs, outcomes), name one failure it sees that the other two do not, from the chapter's
three.
\end{exercise}

\begin{exercise}[$\star$]\label{exo:ml:monitoring-and-retraining:2}
Why is the calibrated threshold on the largest feature index (0.115) above the credit-scoring rule of thumb of 0.10, and what
would a threshold of 0.10 do?
\end{exercise}

\begin{exercise}[$\star$]\label{exo:ml:monitoring-and-retraining:3}
Why does a blind monitor, calibrated to one false page a month, still show a first alarm about two weeks after a failure?
\end{exercise}

\begin{exercise}[$\star\star$]\label{exo:ml:monitoring-and-retraining:4}
Check the chapter's figure of about 150 days for the outcome monitor to see the unit change, from the information coefficients
0.103 and 0.087 and the daily standard deviation 0.071. How many days would a firm trading 2,000 names need?
\end{exercise}

\begin{exercise}[$\star\star$]\label{exo:ml:monitoring-and-retraining:5}
Why does the shadow test give no verdict for ten days? What goes wrong without that rule?
\end{exercise}

\begin{exercise}[$\star\star$]\label{exo:ml:monitoring-and-retraining:6}
\emph{Find the flaw.} ``The PSI alarm went off after the vendor change, so we retrained the model on the new data and the alarm
stopped.''
\end{exercise}

\begin{exercise}[$\star\star\star$]\label{exo:ml:monitoring-and-retraining:7}
\emph{Coding.} Add the pooled monitor of the tutorial's last step (the largest feature index over the last 20 days of samples),
calibrate it to one page a month, and report its delays. Why do its false alarms come in spells?
\end{exercise}

\begin{exercise}[$\star\star\star$]\label{exo:ml:monitoring-and-retraining:8}
Design the monitors for a fade with no change of spread: the third feature's effect falls to zero over 150 days and its
distribution is unchanged. Which layer can see it, how fast, and what would you do instead?
\end{exercise}

\section{Problem: The Tuesday the Model Went Quiet}

\begin{problem}[{Weekend problem --- the Tuesday the model went quiet}]\label{pb:ml:monitoring-and-retraining:1}
The chapter's simulated production run, three planted failures, five monitors.

\medskip\noindent\textbf{Part I --- The model and its monitors.}
\begin{enumerate}
\item Describe the production model and the three failures.
\item What does each layer (inputs, outputs, outcomes) monitor, and when is each available?
\item Define the \omterm{def:ml:monitoring-and-retraining:psi}{population stability index} and give the calibrated thresholds.
\item Why count pages at the onset of an alarm?
\end{enumerate}

\medskip\noindent\textbf{Part II --- \omterm{def:ml:monitoring-and-retraining:performance-alarm}{Performance alarms}.}
\begin{enumerate}[resume]
\item How is the CUSUM on the information coefficient set up and calibrated?
\item How many days does an outcome monitor need to see each failure?
\item Why can outcomes not see the unit change?
\item What is the difference between a first alarm and a sustained alarm?
\end{enumerate}

\medskip\noindent\textbf{Part III --- The delays.}
\begin{enumerate}[resume]
\item Which monitors see the unit change, and when?
\item Which see the reversal?
\item What sees the slow fade, and through what?
\item What would have caught the fade without its change in spread?
\end{enumerate}

\medskip\noindent\textbf{Part IV --- The verdict.}
\begin{enumerate}[resume]
\item State the \emph{named result}: each monitor's detection delay for the three planted failures at one false alarm a month.
\item What happens in shadow after the reversal, and how often is a null challenger promoted?
\item What does a canary save when a faulty challenger is promoted?
\item What should the desk have done on the Tuesday?
\item When should a model be retrained, and when retired?
\item Who should be paged by each monitor?
\item What should the monitoring record in the registry?
\item In one sentence: what does a monitor need before it is worth having?
\end{enumerate}
\end{problem}

\section{Interview questions}

\begin{interviewq}[$\star$ \iqroles{mle}]\label{iq:ml:monitoring-and-retraining:1}
What would you monitor for a model in production, and why at more than one layer?
\end{interviewq}

\begin{interviewq}[$\star\star$ \iqroles{mle, researcher}]\label{iq:ml:monitoring-and-retraining:2}
What is the \omterm{def:ml:monitoring-and-retraining:psi}{population stability index}, and how would you choose its alarm threshold?
\end{interviewq}

\begin{interviewq}[$\star\star$ \iqroles{researcher}]\label{iq:ml:monitoring-and-retraining:3}
Your model's daily IC is 0.10 with a standard deviation of 0.07. How long before you could tell it had fallen to 0.08?
\end{interviewq}

\begin{interviewq}[$\star\star$ \iqroles{mle}]\label{iq:ml:monitoring-and-retraining:4}
Explain \omterm{def:ml:monitoring-and-retraining:shadow}{shadow deployment}, canary and \omterm{def:ml:monitoring-and-retraining:shadow}{champion--challenger}. When is each the right tool?
\end{interviewq}

\begin{interviewq}[$\star\star$ \iqroles{mle, dev}]\label{iq:ml:monitoring-and-retraining:5}
Your input drift alarm fires. Do you retrain?
\end{interviewq}

\begin{interviewq}[$\star\star\star$ \iqroles{mle}]\label{iq:ml:monitoring-and-retraining:6}
Your monitors page the team twenty times a week and nobody looks any more. What do you change?
\end{interviewq}
